\article{（拟）分裂群的“广义 Ramanujan 猜想”的一个反例}{A Counterexample to the ``Generalized Ramanujan Conjecture'' for (Quasi-) Split Groups}{R. Howe 与 I. I. Piatetski-Shapiro}{1}{315--322}

\origsection{1}{引言（Introduction）}
在 [Sat] 中，Satake 说明了如何用群论语言表述关于尖点模形式（cuspidal modular form）系数的 Ramanujan 猜想和 Ramanujan--Petersson 猜想。这一表述的核心是：与经典尖点形式相联系的自守表示（automorphic representation）的局部分量（local constituent；见 [F]）应当是 Harish-Chandra 意义下的缓增表示（tempered）[KZ]（亦见 [GGP]）。缓增性是关于矩阵系数（matrix coefficient）渐近行为的一个技术条件，事实证明，它在 Harish-Chandra 关于 Plancherel 公式的研究中至关重要。由于这一条件对任意约化群（reductive group）都有意义，人们便将“整体域上的任意代数群（global algebraic group），其尖点自守表示是否都缓增”这一问题称为广义 Ramanujan 猜想。若干作者已构造出在无穷处的分量不缓增的自守尖点形式（见 [Ga]、[HW]、[M] 以及 [B] 中的评述）。然而，这些例子属于各向异性群（anisotropic group），拟分裂群（quasi-split group）的情形则仍未解决。本文将构造 $\operatorname{Sp}_4$ 上几乎处处不缓增的尖点形式。\footnote{我们最初的例子是针对 $U_{2,1}$ 的；不过，$\operatorname{Sp}_4$ 的例子更容易陈述，而且在某种意义上更好，因为 $\operatorname{Sp}_4$ 不仅拟分裂，而且分裂（split）。}

我们的构造使用振子表示（oscillator representation）。把 $\operatorname{Sp}_4$ 放入约化对偶对（dual reductive pair）$(\operatorname{Sp}_4,O_2)$ 中（见 [H]），即可构造从 $O_2$ 上的自守形式到 $\operatorname{Sp}_4$ 上的自守形式的一个单射。这一方法与 Hecke、Maass 及 Shalika--Tanaka 的构造相呼应：他们从二次域的广义 Hecke 特征（Grössencharakter）出发，构造 $\operatorname{Sp}_2=\operatorname{SL}_2$ 上的尖点形式。事实上，后者正是与对偶对 $(\operatorname{Sp}_2,O_2)$ 对应的 $\operatorname{SL}_2$ 上的自守形式。我们指出，虽然 $O_2$ 的每个自守表示都给出 $\operatorname{Sp}_4$ 的一个表示，但只有其中一部分出现在与 $\operatorname{Sp}_2$ 的对应中。结果表明，在与 $\operatorname{Sp}_2$ 的对应中没有出现的 $O_2$ 表示，恰好产生 $\operatorname{Sp}_4$ 的尖点表示。局部地说，在二元二次型为各向异性的一个素位（place），$O_2$ 恰有一个表示没有出现在与 $\operatorname{Sp}_2$ 的配对中，即在 $\operatorname{SO}_2$ 上平凡的非平凡特征（character）。因此，$O_2$ 的非连通性是这一构造的关键。C. Asmuth 在其学位论文 [A] 中分析了局部的 $(\operatorname{Sp}_4,O_2)$ 对应。

所构造的尖点形式有若干特殊之处。特别地，它们没有 Whittaker 模型（Whittaker model）[P1]。因而，对于具有 Whittaker 模型的尖点自守表示，广义 Ramanujan 猜想仍有可能成立。没有 Whittaker 模型的尖点表示中，最著名的例子之一是 Srinivasan [Sr] 发现的有限域上 $\operatorname{Sp}_4$ 的表示 $\theta_{10}$。这里构造的尖点形式与 $\theta_{10}$ 密切相关。事实上，设 $k$ 是非阿基米德局部域（non-Archimedean local field），$\mathcal O$ 为其整数环，$F$ 为其剩余域（residue class field）。那么 $\theta_{10}$ 是 $\operatorname{Sp}_4(F)$ 的一个表示，可以经由典范满射
\[
\operatorname{Sp}_4(\mathcal O)\longrightarrow\operatorname{Sp}_4(F)
\]
拉回为 $\operatorname{Sp}_4(\mathcal O)$ 的表示。由 Cartan 分解（Cartan decomposition）可知诱导表示（induced representation）
\[
\operatorname{ind}\theta_{10}
 =\operatorname{ind}_{\operatorname{Sp}_4(\mathcal O)}^{\operatorname{Sp}_4(k)}\theta_{10}
\]
是不可约（irreducible）且超尖点（supercuspidal）的。在我们构造的自守表示中，某些有限素位的局部分量本质上就是这个 $\operatorname{ind}\theta_{10}$。此外，$\theta_{10}$ 是 Lusztig [L] 称为幺幂（unipotent）的那类表示中最简单的尖点表示之一。$U_{2,1}$ 的例子也涉及该群的一个尖点幺幂表示。这提示我们，广义 Ramanujan 猜想的失效与这些尖点幺幂表示的存在有关。Lusztig 已对有限域上经典群的这些表示作出分类，并证明，虽然这类表示并不常见，但在秩任意大的某些群中仍然存在。

另一个值得研究的问题，是计算这些自守形式的 $L$ 函数。使用 [P2]、[NPS] 中对 $\operatorname{GSp}_4$ 的 $L$ 函数的定义，我们做了一些初步计算。这些计算提示，来自振子表示的局部尖点表示的 $L$ 函数有极点。这与如下预期相符：经过适当提升（lifting）后，我们的尖点但不缓增的自守表示对应于一个非尖点自守形式。这也可能说明，本文的例子如何融入 Langlands 理论的整体构想。详细计算将另文给出。

\origsection{2}{构造（The construction）}
设 $k$ 为整体域（global field），即 Weil [W2] 所谓的 $A$ 域。假设 $\operatorname{char}k\ne2$。令 $\operatorname{Sp}_4(k)$ 为元素在 $k$ 中且形如
\[
\begin{pmatrix}A&B\\C&D\end{pmatrix}
\]
的矩阵全体，其中 $A,B,C,D$ 是满足
\begin{equation}
 A^tD-C^tD=I,\qquad A^tC=C^tA,\qquad B^tD=D^tB
 \tag{2.1}
\end{equation}
的 $2\times2$ 矩阵。这里 $I$ 是 $2\times2$ 单位矩阵，$A^t$ 表示 $A$ 的转置。\footnote{译注：原书式 (2.1) 的第一式印作 $A^tD-C^tD=I$；本译文保留原式。通常辛群的该条件写作 $A^tD-C^tB=I$。}

设 $K$ 是 $k$ 的二次扩张。从 $K$ 到 $k$ 的范数型（norm form）是一个二次型，记作 $\beta$，这里将 $K$ 视为 $k$ 上的二维向量空间。将这个型的等距群（isometry group）记作 $O_2(k)$。令 $\tau$ 为 $K/k$ 的 Galois 群 $\operatorname{Gal}$ 的非平凡元素。那么 $\tau\in O_2(k)$，并且对 $x\in K$ 有 $\beta(x,x)=x\tau(x)$。令
\begin{equation}
 U=\{x\in K:\beta(x,x)=1\}.\tag{2.2}
\end{equation}
显然，$U$ 是 $K$ 的乘法群 $K^\times$ 的一个子群，因而通过乘法作用于 $K$。这一作用给出 $U$ 与 $\operatorname{SO}_2$ 的同构，并且有正合列（exact sequence）
\begin{equation}
 1\longrightarrow U\longrightarrow O_2\longrightarrow\operatorname{Gal}\longrightarrow1.
 \tag{2.3}
\end{equation}
令 $A$ 为 $k$ 的阿代尔环（adele ring）。用 $v$ 表示 $k$ 的任意素位，用 $k_v$ 表示 $k$ 在 $v$ 处的完备化（completion）。那么 $\operatorname{Sp}_4(A)$ 是由元素在 $A$ 中且满足 (2.1) 的矩阵组成的群，$K_v=K\otimes_k k_v$ 是由 $K$ 扩张标量得到的二次 $k_v$ 代数，$\beta_v$ 是 $\beta$ 在 $K_v$ 上诱导的型，其余记号依此类推。

根据 Weil [W1] 和 Saito [So] 的工作，我们知道，$\operatorname{Sp}_4(A)$ 在
\[
 L^2(A^2\otimes_k K)\simeq L^2(A^2\otimes_A K_A)\simeq L^2(K_A^2)
\]
上有一个酉表示（unitary representation）$\omega$。由上述工作可知，这个表示具有以下性质。

\begin{enumerate}
\item[(i)] $\omega$ 保持 Schwartz--Bruhat 空间（Schwartz--Bruhat space）$\mathcal S(K_A^2)$。以 $\omega^\infty$ 表示 $\omega$ 在 $\mathcal S$ 上的限制。
\item[(ii)] 正如 $\mathcal S(K_A^2)$ 是局部 Schwartz--Bruhat 空间 $\mathcal S(K_v^2)$ 的限制张量积（restricted tensor product），$\omega^\infty$ 也是 $\operatorname{Sp}_4(k_v)$ 作用在 $\mathcal S(K_v^2)$ 上的表示 $\omega_v^\infty$ 的限制张量积；而 $\omega_v^\infty$ 是 $L^2(K_v^2)$ 上的酉表示 $\omega_v$ 的限制。
\item[(iii)] $\operatorname{Sp}_4$ 中形如
\[
 \begin{pmatrix}A&0\\0&A^{t-1}\end{pmatrix}
\]
的元素组成的子群，其作用由 $A^2$ 上的线性作用诱导，再乘以适当的标量以保证酉性。具体来说，选择 $K/k$ 的一组基 $e_1,e_2$。那么，对 $f\in\mathcal S(A^2\otimes_A K_A)$，有
\begin{equation}
\begin{split}
 &\omega\!\begin{pmatrix}A&0\\0&A^{t-1}\end{pmatrix}(f)
       (v_1\otimes e_1+v_2\otimes e_2)\\
 &\qquad=s(A)|\det A|^{-1}
 f\bigl(A^{-1}(v_1)\otimes e_1+A^{-1}(v_2)\otimes e_2\bigr).
\end{split}\tag{2.4}
\end{equation}
在 (2.4) 中，$|\ |$ 是 $A$ 上标准的“模”（module；参见 [W2]）或绝对值，$v_1,v_2\in A^2$，而 $s(A)$ 是 $\operatorname{GL}_2(A)$ 上一个二阶特征。式 (2.4) 自然分解为各素位处的局部作用。
\item[(iv)] 由 $O_2(A)$ 在 $K_A$ 上的线性作用诱导的 $O_2(A)$ 在 $\mathcal S(K_A^2)$ 上的作用 $\rho$，与 $\omega(\operatorname{Sp}_4(A))$ 对易。注意，$O_2(A)$ 的这个作用同样分解为局部 $O_2(k_v)$ 的作用之积。
\item[(v)] 对 $f\in\mathcal S(K_A^2)$，定义
\begin{equation}
 \Theta(f)=\sum_{x\in K^2}f(x).\tag{2.5}
\end{equation}
那么 $\Theta$ 是 $\mathcal S(K_A^2)$ 上的一个分布（distribution），并且在 $\omega^\infty(\operatorname{Sp}_4(k))$ 下不变。也就是说，对 $g\in\operatorname{Sp}_4(k)$ 和 $f\in\mathcal S(K_A^2)$，有
\begin{equation}
 \Theta(\omega(g)(f))=\Theta(f).\tag{2.6}
\end{equation}
\end{enumerate}

显然，$\Theta$ 在 $\rho(O_2(k))$ 下也不变。因此，如果对 $f\in\mathcal S(K_A^2)$，由公式
\begin{equation}
 \theta_f(g,h)=\Theta(\omega(g)\rho(h)f)\tag{2.7}
\end{equation}
定义 $\operatorname{Sp}_4(A)\times O_2(A)$ 上的函数 $\theta_f$，那么 $\theta_f$ 下降为商空间
\[
 \bigl(\operatorname{Sp}_4(k)\times O_2(k)\bigr)
 \backslash\bigl(\operatorname{Sp}_4(A)\times O_2(A)\bigr)
\]
上的函数。此外，$\theta_f$ 在无穷处具有适度增长（moderate growth），所以它是 $\operatorname{Sp}_4(A)\times O_2(A)$ 上的自守形式。\footnote{\textbf{译者注（本次校订）：}对任意 Schwartz--Bruhat 函数 $f$，这里首先得到的是光滑、适度增长的自守函数。若按本书通常要求的 $K$ 有限性与包络代数中心有限性来定义“自守形式”，还须选择相应的有限向量或进行谱投影；适度增长本身不能代替这两个条件。}

下面采用 $\theta$ 级数（theta-series）理论中的标准构造。陪集空间 $O_2(k)\backslash O_2(A)$ 是紧的。下文将证明，适当选择 $f$，便可使$\theta_f$ 在 $O_2(A)$ 上的限制成为 $O_2(k)\backslash O_2(A)$ 上任意的 Schwartz--Bruhat 函数。特别地，可以使 $\theta_f$ 按照 $O_2(A)$ 的某个给定不可约表示 $\sigma$ 变换。研究局部表示 $\omega_v$ 后，即可推出：若在某些素位 $v$ 适当选择 $\sigma_v$，则 $\theta_f$ 在 $\operatorname{Sp}_4(k_v)$ 下按照某个超尖点表示变换。\footnote{实际上，可以证明 $\theta_f$ 在 $\operatorname{Sp}_4(A)$ 下按照由 $\sigma$ 决定的一个不可约表示 $\sigma'$ 变换；不过，这个更精细的结论并非本文所需，证明也会占用较多篇幅。} 因而，$\theta_f|\operatorname{Sp}_4(A)$ 是尖点形式。进一步分析表明，在几乎所有素位 $v$，尖点形式 $\theta_f$ 都按照由 $\sigma_v$ 决定的 $\operatorname{Sp}_4(k_v)$ 的某个不可约而不缓增的表示 $\sigma_v'$ 变换。所以，在 $\operatorname{Sp}_4(A)$ 上由 $\theta_f$ 生成的尖点形式空间中，任意不可约分量在几乎所有素位都按照这些相同的 $\sigma_v'$ 变换，这就得到所需的反例。下面给出证明。

\begin{originaltheorem}{命题 1（Proposition 1）}
映射 $f\mapsto\theta_f|O_2(A)$ 从 $\mathcal S(K_A^2)$ 满射到 $\mathcal S(O_2(k)\backslash O_2(A))$。
\end{originaltheorem}
\begin{proof}
取 $K^2$ 中一点 $(x_1,x_2)$，使 $x_1,x_2$ 在 $k$ 上张成 $K$。从 $O_2(A)$ 到 $K_A^2$ 的映射
\[
 e:g\longmapsto(gx_1,gx_2)
\]
是光滑单射，并且其像是闭的。这一事实可由以下两点看出：$O_2(k)$ 的像由有理点构成，因而离散，而 $O_2(k)\backslash O_2(A)$ 是紧的。由于 $e$ 的像是闭的，拉回
\[
 e^*:C_c^\infty(K_A^2)\longrightarrow C^\infty(O_2(A))
\]
满射到 $C_c^\infty(O_2(A))$。取 $f\in C_c^\infty(K_A^2)$，使其在不属于 $e$ 的像的所有 $K^2$ 中的点处均为零。由 Witt 定理 [J] 可得
\begin{equation}
 \theta_f(g)=\sum_{h\in O_2(k)}e^*(f)(hg).\tag{2.8}
\end{equation}
$C^\infty(O_2(k)\backslash O_2(A))$ 中的任意函数都可写成 (2.8) 右端的周期化形式，因此命题成立。
\end{proof}

\paragraph*{注（Remark）.}
对于任意不可约的 I 型约化对偶对 $(G,G')$（见 [H]），只要与 $G$ 相联系的带型空间（space with form）具有一个各向同性子空间（isotropic subspace），且其维数至少等于与 $G'$ 对应的带型空间的维数，上述结论及证明便同样适用。

如果 $K_v$ 是域，也就是说，$K$ 在 $v$ 处不分裂，那么 $\beta_v$ 是各向异性的，$O_2(k_v)$ 是紧群。因此，$\mathcal S(K_v^2)$ 按 $O_2(k_v)$ 的不可约表示分解为离散直和；每个等型分量对应于该群的一个一维或二维不可约表示。Asmuth 已证明以下事实 [A]（亦见 [H]）。

\begin{originaltheorem}{命题 2（Proposition 2）}
如果 $K/k$ 在 $v$ 处不分裂，则有正交直和分解\footnote{\textbf{译者注（本次校订）：}本命题中关于“超尖点表示”的陈述，以及随后利用紧支集与局部常值性化简 (2.9) 的论证，须将 $v$ 限定为非阿基米德素位。整体反例的构造可选择这样的非分裂有限素位。}
\[
 \mathcal S(K_v^2)\simeq\sum_\sigma\mathcal S(K_v^2,\sigma),
\]
其中 $\sigma$ 遍历 $O_2(k_v)$ 的所有不可约表示。每个空间 $\mathcal S(K_v^2,\sigma)$ 在 $O_2(k_v)$ 与 $\operatorname{Sp}_4(k_v)$ 的联合作用下不可约（后者通过 $\omega_v$ 在 $\mathcal S(K_v^2,\sigma)$ 上的限制作用），因而作为 $O_2\times\operatorname{Sp}_4$ 模具有 $\sigma\otimes\sigma'$ 的形式。映射 $\sigma\mapsto\sigma'$ 定义了从 $O_2(k_v)$ 的表示到 $\operatorname{Sp}_4(k_v)$ 的（不可约、容许（admissible）、可酉化（unitarizable））表示的一个单射。如果 $\sigma$ 是 $O_2$ 的一个非平凡一维表示，且在 $\operatorname{SO}_2$ 上平凡（即符号表示（signum representation）），那么 $\sigma'$ 是超尖点表示。
\end{originaltheorem}

若 $\sigma$ 是 $O_2$ 的任意一个不同于符号表示的表示，则相应的 $\sigma'$ 甚至不缓增 [KZ]。下面只证明 $\sigma'$ 为平凡表示时的结论，这已足以构造所需反例。\footnote{译注：此处原文为“when $\sigma'$ is the trivial representation”；下文实际选择的是 $O_2(k_v)$ 不变函数，因而所讨论的是 $\sigma$ 为平凡表示的情形。这里保留原文表述并指出这一记号疑误。}关键是选取函数 $f,f'\in\mathcal S(K_v^2)$，使其在 $O_2(k_v)$ 下不变、取非负值，并在原点非零。计算矩阵系数，得到
\begin{equation}
\begin{split}
 &\left\langle f,\omega_v\!\begin{pmatrix}
 r&0&0&0\\0&t&0&0\\0&0&r^{-1}&0\\0&0&0&t^{-1}
 \end{pmatrix}(f')\right\rangle\\
 &\qquad=\pm|rt|_v^{-1}\int_{K_v^2}
 f(x_1,x_2)\overline{f'(r^{-1}x_1,t^{-1}x_2)}\,dx_1\,dx_2.
\end{split}\tag{2.9}
\end{equation}
由于 $f$ 有紧支集（compact support），当 $|r|_v,|t|_v$ 充分大时，(2.9) 变为
\[
 \pm|rt|_v^{-1}\overline{f'(0,0)}
 \int_{K_v^2}f(x_1,x_2)\,dx_1\,dx_2.
\]
由于 $\operatorname{Sp}_4$ 的 $\delta$ 函数是 $|r^4t^2|_v$，可见矩阵系数 (2.9) 在无穷处衰减得太慢，因此 $\sigma'$ 不可能缓增。

如果 $K_v$ 不是域，那么 $\beta_v$ 是双曲平面（hyperbolic plane），其两条各向同性直线分别对应于 $K$ 中位于 $v$ 之上的两个素位。$O_2$ 与 $\operatorname{Sp}_4$ 在 $\mathcal S(K_v^2)$ 上的联合作用可以用与命题 2 完全类似的方式分析，此时须使用连续分解，论证因而稍长。由于缺少可直接引用的文献，下面证明所需的部分，并只考虑非阿基米德素位。令 $C,C'$ 分别为 $O_2(k_v)$ 和 $\operatorname{Sp}_4(k_v)$ 中的整矩阵子群。以 $\mathcal S(K_v^2)^{C'}$ 表示 $\mathcal S(K_v^2)$ 中的 $C'$ 不变向量。

\begin{originaltheorem}{命题 3（Proposition 3）}
存在一个线性同构
\[
 \alpha:C_c^\infty(O_2(k_v)/C)\longrightarrow\mathcal S(K_v^2)^{C'}.
\]
映射 $\alpha$ 是一个 $O_2(k_v)$ 交织映射（intertwining map）。
\end{originaltheorem}
\begin{proof}
进行部分 Fourier 变换（partial Fourier transform；$(O_2,\operatorname{SL}_2)$ 对的相似情形见 [RS]），可知 $\omega_v^\infty$ 等价于由 $\operatorname{Sp}_4(k_v)$ 在 $k_v^4$ 上的自然线性作用所诱导的、该群在 $\mathcal S(k_v^4)$ 上的表示。在 $\omega_v^\infty$ 的这一实现中，$\operatorname{SO}_2(k_v)$ 的作用 $\rho_v^\infty$ 变为 $k_v^4$ 的（标量）伸缩，同时乘以保证酉性的归一化因子。再加入以下算子，便得到整个 $O_2$ 的作用： $k_v^4$ 上相对于辛型（symplectic form）
\[
 \begin{pmatrix}0&I\\-I&0\end{pmatrix}
\]
的 Fourier 变换；$\operatorname{Sp}_4(k_v)$ 正是该型的等距群。

容易验证，$C'$ 在 $k_v^4$ 中的轨道是集合 $\pi^k(\mathcal O_v^4-\pi\mathcal O_v^4)$，以及 $0$。由此可见，$\mathcal S(k_v^4)^{C'}$ 有一组由函数 $\phi_k$ 构成的基，其中 $\phi_k$ 是 $\pi^k\mathcal O_v^4$ 的特征函数。所有 $\phi_k$ 都在 $\mathcal O_v^\times$ 对 $k_v^4$ 的标量伸缩作用下不变。适当归一化 $k_v^4$ 上的 Fourier 变换 $\widehat{\ }$ 后，有
\begin{equation}
 \widehat\phi_k=|\pi|_v^{4k}\phi_{-k}.\tag{2.10}
\end{equation}
特别地，$\phi_0$ 在 $\mathcal O_v^\times$ 和 $\widehat{\ }$ 下均不变，也就是在 $C$ 下不变。定义映射
\[
 \alpha:C_c^\infty(O_2(k_v)/C)\longrightarrow\mathcal S(k_v^4)^{C'}
\]
为
\[
 \alpha(f)=\rho_v(f)(\phi_0),\qquad f\in C_c^\infty(O_2(k_v)/C).
\]
由于 $\pi$ 的幂所给出的伸缩构成 $\rho_v(O_2(k_v))/\rho_v(C)$ 的一组陪集代表，容易看出 $\alpha$ 是线性同构。它显然也是一个 $O_2(k_v)$ 交织映射，命题得证。
\end{proof}

若 $\operatorname{Sp}_4(k_v)$ 的不可约容许表示包含非零的 $C'$ 不变向量，就称它为 $C'$ 球表示（$C'$-spherical representation）。对 $O_2(k_v)$ 也使用相似的术语。

\begin{originaltheorem}{命题 4（Proposition 4）}
设 $\sigma'$ 是 $\operatorname{Sp}_4(k_v)$ 的一个（不可约容许的）$C'$ 球表示。\footnote{\textbf{译者注（本次校订）：}这里还须假设 $\operatorname{Hom}_{\operatorname{Sp}_4(k_v)}(\mathcal S(K_v^2),\sigma')\ne0$，即 $\sigma'$ 确实出现在这一对应中。否则所有交织算子组成的空间为零，其共同核是整个 $\mathcal S(K_v^2)$，从而 $Q_{\sigma'}=0$，不能称为不可约表示。后面的 (2.12) 也表明，能出现的球表示的两个 Hecke 特征值须满足特定关系。}令 $Y_{\sigma'}$ 为从 $\mathcal S(K_v^2)$ 到 $\sigma'$ 的空间的所有 $\operatorname{Sp}_4(k_v)$ 交织算子（intertwining operator）的共同核。那么，$O_2(k_v)\times\operatorname{Sp}_4(k_v)$ 在
\[
 Q_{\sigma'}=\mathcal S(K_v^2)/Y_{\sigma'}
\]
上的联合作用不可约，且具有 $\sigma\otimes\sigma'$ 的形式，其中 $\sigma$ 是 $O_2(k_v)$ 的一个不可约 $C$ 球表示。对应 $\sigma\mapsto\sigma'$ 是从 $O_2(k_v)$ 的所有 $C$ 球表示的集合到 $\operatorname{Sp}_4(k_v)$ 的 $C'$ 球表示的集合的一个单射。
\end{originaltheorem}
\begin{proof}
由于 $C'$ 紧，自然投影 $\mathcal S(K_v^2)^{C'}\to Q_{\sigma'}^{C'}$ 是满射。空间 $Q_{\sigma'}^{C'}$ 在 $O_2(k_v)$ 下不变；又因为 $\sigma'$ 的 $C'$ 不变向量空间是一维的 [C]，$Q_{\sigma'}$ 在 $O_2\times\operatorname{Sp}_4$ 下不可约，当且仅当 $Q_{\sigma'}^{C'}$ 在 $O_2$ 下不可约。由命题 3，有满射的 $O_2$ 交织映射
\[
 \alpha':C_c^\infty(O_2/C)\longrightarrow Q_{\sigma'}^{C'}.
\]
因而，$Q_{\sigma'}^{C'}$ 中的不可约 $O_2$ 表示与 $Q_{\sigma'}^{C'}$ 中的 $C$ 不变向量对应。所以，要证明 $Q_{\sigma'}$ 在 $O_2\times\operatorname{Sp}_4$ 下不可约，只须证明 $Q_{\sigma'}^{C'\times C}$ 是一维的。为此，考虑 Hecke 代数（Hecke algebra）[C] $C_c^\infty(O_2//C)$ 和 $C_c^\infty(\operatorname{Sp}_4//C')$ 在命题 3 所描述的 $\mathcal S(k_v^4)^{C\times C'}$ 上的作用。（下文交替使用 $\mathcal S(k_v^4)$ 与 $\mathcal S(K_v^2)$ 这两个等价模型。）若 $\phi_k$ 如 (2.10) 所定义，则容易看出，函数
\[
 q^{2k}\phi_k+q^{-2k}\phi_{-k},\qquad q=|\pi|_v^{-1},
\]
在 $k\ge0$ 时构成 $\mathcal S(k_v^4)^{C\times C'}$ 的一组基。令 $T_1$ 为 $O_2$ 上如下定义的 Hecke 算子：
\[
 T_1(f)(x)=q^{-2}f(\pi x)+q^2f(\pi^{-1}x),
 \qquad f\in\mathcal S(k_v^4)^{C\times C'}.
\]
那么 $T_1$ 生成 Hecke 代数 $C_c^\infty(O_2//C)$。令 $T_1',T_2'$ 为 $\operatorname{Sp}_4$ 上的 Hecke 算子，各自的总质量为 $1$，支集分别为以下代表元所对应的 $C'$ 双陪集：
\[
 \begin{pmatrix}\pi&0&0&0\\0&1&0&0\\0&0&\pi^{-1}&0\\0&0&0&1\end{pmatrix},
 \qquad
 \begin{pmatrix}\pi&0&0&0\\0&\pi&0&0\\0&0&\pi^{-1}&0\\0&0&0&\pi^{-1}\end{pmatrix}.
\]
这些算子生成 $C_c^\infty(\operatorname{Sp}_4//C')$。

直接计算得到
\begin{equation}
\begin{aligned}
 T_1(\phi_0)&=q^{-2}\phi_{-1}+q^2\phi_1,\\
 T_1'(\phi_0)&=\frac{(q-1)\phi_{-1}+(q^4-2q+1)\phi_0+q^4(q-1)\phi_1}{q(q^4-1)},\\
 T_2'(\phi_0)&=\frac{(q^2-1)\phi_{-1}+(q^4-2q^2+1)\phi_0+q^4(q^2-1)\phi_1}{q^2(q^4-1)}.
\end{aligned}\tag{2.11}
\end{equation}
由命题 3 可知，$\mathcal S(k_v^4)^{C\times C'}$ 作为 $C_c^\infty(O_2//C)$ 模，同构于 $C_c^\infty(O_2//C)$ 在自身上的正则作用。特别地，在 $T_1$ 的作用下，$\phi_0$ 是 $\mathcal S(k_v^4)^{C\times C'}$ 的循环向量（cyclic vector）。由于 $T_1$ 与各 $T_j'$ 两两对易，对任意 $f\in\mathcal S(k_v^4)^{C\times C'}$，有
\begin{equation}
\begin{aligned}
 q(q^4-1)T_1'(f)&=(q^4-2q+1)f+q^2(q-1)T_1(f),\\
 q^2(q^4-1)T_2'(f)&=(q^4-2q^2+1)f+q^2(q^2-1)T_1(f).
\end{aligned}\tag{2.12}
\end{equation}
再考察 $Q_{\sigma'}^{C\times C'}$。它是 $\mathcal S(K_v^2)^{C\times C'}$ 的商。因此，若 $\bar\phi_0$ 为 $\phi_0$ 在 $Q_{\sigma'}^{C\times C'}$ 中的像，那么 $\bar\phi_0$ 是 $T_1$ 作用在 $Q_{\sigma'}^{C\times C'}$ 上的循环向量。另一方面，由于 $Q_{\sigma'}$ 作为 $\operatorname{Sp}_4$ 模是若干个 $\sigma'$ 的直和，可见 $\bar\phi_0$ 是 $T_1'$ 和 $T_2'$ 的特征向量（eigenvector）。于是，由 (2.11)，$\bar\phi_0$ 也是 $T_1$ 的特征向量，所以 $Q_{\sigma'}^{C\times C'}$ 确为一维。因此 $Q_{\sigma'}\simeq\sigma\otimes\sigma'$。此外，(2.12) 表明 $T_1$ 在 $\bar\phi_0$ 上的特征值决定 $T_1'$ 和 $T_2'$ 的特征值，反之亦然。由于这些特征值又分别决定 $\sigma$ 和 $\sigma'$ [C]，可知 $\sigma\mapsto\sigma'$ 也是单射，命题得证。
\end{proof}

由 $\omega_v^\infty$ 在 $\mathcal S(k_v^4)$ 上的实现还可看出，它的任何 $C'$ 球商都不缓增，因为所有这类商都容许一个在 $\operatorname{Sp}_4$ 中 $k_v^4$ 的某一点的稳定子群（isotropy group）下不变的分布，而这种分布的存在排除了缓增性。

现在可以完成整体构造。选择 $f\in\mathcal S(K_A^2)$，使 $\theta_f$ 在 $O_2(A)$ 下按照一个不可约表示变换，并且这个表示在至少一个 $K/k$ 不分裂的素位上的局部分量是符号表示。这样的表示显然存在。由命题 2 和命题 4，可知在所有 $K/k$ 不分裂的素位 $v$，或者在 $K$ 分裂且 $\theta_f$ 被 $C_v'$ 固定的素位，$\theta_f$ 在 $\operatorname{Sp}_4(A)$ 下都按照由 $\sigma_v$ 决定的 $\operatorname{Sp}_4(k_v)$ 的不可约表示 $\sigma_v'$ 变换。这些素位已包含几乎所有素位。由于 $\theta_f$ 在至少一个素位按照超尖点表示变换，它必为尖点形式。因此，$\theta_f$ 的所有 $\operatorname{Sp}_4(A)$ 平移所成空间的 $L^2$ 闭包，分解为尖点自守表示的直和。在上述所有素位，这些自守表示的局部分量必为由 $\sigma_v$ 决定的 $\sigma_v'$。因此，所得自守表示几乎处处不缓增。

最后指出，若 $k=\Q$，且 $K$ 是虚二次扩张，那么我们构造的某些形式就是权为 $1$ 的经典 Siegel 模形式（Siegel modular form）。

此外，由于 $O_2$ 上的自守形式显然不满足强重数一性质（strong multiplicity one），相应的 $\operatorname{Sp}_4$ 上的形式也违反强重数一性质。

\begin{thebibliography}{GGP}
\bibitem[A]{c20-A} C. Asmuth, \emph{Weil representations of symplectic $p$-adic groups}, 预印本。
\bibitem[B]{c20-B} A. Borel, \emph{Automorphic $L$-functions}, 本论文集，第 2 册，27--61 页。
\bibitem[C]{c20-C} P. Cartier, \emph{Representations of $p$-adic groups: A survey}, 本论文集，第 1 册，111--155 页。
\bibitem[F]{c20-F} D. Flath, \emph{Decomposition of representations into tensor products}, 本论文集，第 1 册，179--183 页。
\bibitem[Ga]{c20-Ga} R. Gangolli, \emph{Zeta-functions of Selberg's type for compact space forms of symmetric spaces of rank one}, Illinois J. Math. \textbf{21} (1977), 1--41.
\bibitem[GGP]{c20-GGP} I. M. Gelfand, M. I. Graev and I. Piatetski-Shapiro, \emph{Representation theory and automorphic functions}, W. B. Saunders Co., 1969.
\bibitem[HW]{c20-HW} R. Hotta and N. Wallach, \emph{On Matsushima formula for the Betti numbers of a locally symmetric space}, Osaka J. Math.
\bibitem[H]{c20-H} R. Howe, \emph{$\theta$-series and invariant theory}, 本论文集，第 1 册，275--285 页。
\bibitem[J]{c20-J} N. Jacobson, \emph{Lectures in abstract algebra. II, Linear Algebra}, Springer-Verlag, Berlin, 1975.
\bibitem[KZ]{c20-KZ} A. W. Knapp and G. Zuckerman, \emph{Normalizing factors, tempered representations and $L$-groups}, 本论文集，第 1 册，93--105 页。
\bibitem[L]{c20-L} G. Lusztig, \emph{Irreducible representations of finite classical groups}, 预印本。
\bibitem[M]{c20-M} J. Millson, \emph{On the first Betti number of a constant negatively curved manifold}, Ann. of Math. (2) \textbf{104} (1976), 235--247.
\bibitem[NPS]{c20-NPS} M. Novodvorsky and I. Piatetski-Shapiro, \emph{Rankin-Selberg method in the theory of automorphic forms}, Proc. Sympos. Pure Math., vol. 30, Amer. Math. Soc., Providence, R.I., 1977.
\bibitem[P1]{c20-P1} I. Piatetski-Shapiro, \emph{Multiplicity one theorems}, 本论文集，第 1 册，209--212 页。
\bibitem[P2]{c20-P2} I. Piatetski-Shapiro, \emph{Euler subgroup}, Lie groups and their representations, Summer School in Budapest (1971).
\bibitem[RS]{c20-RS} S. Rallis and G. Schiffman, \emph{Weil Representation. I, Intertwining distributions and discrete spectrum}, 预印本。
\bibitem[So]{c20-So} M. Saito, \emph{Représentations unitaires des groupes symplectiques}, J. Math. Soc. Japan \textbf{24} (1972), 232--251.
\bibitem[Sat]{c20-Sat} I. Satake, \emph{Spherical functions and Ramanujan conjecture}, Proc. Sympos. Pure Math., vol. 9, Amer. Math. Soc., Providence, R.I., 1966, pp. 258--264.
\bibitem[Sr]{c20-Sr} B. Srinivasan, \emph{The characters of the finite symplectic group $\operatorname{Sp}(4,q)$}, Trans. Amer. Math. Soc. \textbf{131} (1968), 488--525.
\bibitem[W1]{c20-W1} A. Weil, \emph{Sur certain groupes d'opérateurs unitaires}, Acta Math. \textbf{111} (1964), 143--211.
\bibitem[W2]{c20-W2} A. Weil, \emph{Basic number theory}, Springer-Verlag, Berlin, 1973.
\end{thebibliography}

\noindent Yale 大学（Yale University）\\
Tel Aviv 大学（University of Tel Aviv）
